\documentclass[a4paper,12pt]{article}

% =============================================================================
% Pacotes
% =============================================================================
\usepackage[utf8]{inputenc}
\usepackage[T1]{fontenc}
\usepackage[brazilian,provide=*]{babel}
\usepackage[margin=2.5cm]{geometry}
\usepackage{graphicx}
\usepackage{listings}
\usepackage{xcolor}
\usepackage{hyperref}
\usepackage{titlesec}
\usepackage{indentfirst}
\usepackage{enumitem}

% =============================================================================
% Configuração de links e cores
% =============================================================================
\hypersetup{
    colorlinks=true,
    linkcolor=black,
    citecolor=black,
    urlcolor=blue,
    pdftitle={Linux Mint em Docker - Estrutura do Sistema Operacional e Kernel},
    pdfauthor={Equipe AV2}
}

% =============================================================================
% Configuração de blocos de código (Shell Script)
% =============================================================================
\lstdefinestyle{shell}{
    backgroundcolor=\color{black!5},
    basicstyle=\ttfamily\footnotesize,
    commentstyle=\color{green!40!black},
    keywordstyle=\color{blue},
    stringstyle=\color{red!60!black},
    numbers=left,
    numberstyle=\tiny\color{gray},
    breaklines=true,
    frame=single,
    tabsize=2,
    language=bash,
    literate=
        {á}{{\'a}}1 {à}{{\`a}}1 {ã}{{\~a}}1 {â}{{\^a}}1 {ä}{{\"a}}1
        {é}{{\'e}}1 {è}{{\`e}}1 {ê}{{\^e}}1 {ë}{{\"e}}1
        {í}{{\'i}}1 {ì}{{\`i}}1 {î}{{\^i}}1 {ï}{{\"i}}1
        {ó}{{\'o}}1 {ò}{{\`o}}1 {õ}{{\~o}}1 {ô}{{\^o}}1 {ö}{{\"o}}1
        {ú}{{\'u}}1 {ù}{{\`u}}1 {û}{{\^u}}1 {ü}{{\"u}}1
        {ç}{{\c c}}1
        {Á}{{\'A}}1 {À}{{\`A}}1 {Ã}{{\~A}}1 {Â}{{\^A}}1
        {É}{{\'E}}1 {È}{{\`E}}1 {Ê}{{\^E}}1
        {Í}{{\'I}}1 {Ì}{{\`I}}1 {Î}{{\^I}}1
        {Ó}{{\'O}}1 {Ò}{{\`O}}1 {Õ}{{\~O}}1 {Ô}{{\^O}}1
        {Ú}{{\'U}}1 {Ù}{{\`U}}1 {Û}{{\^U}}1
        {Ç}{{\c C}}1
        {ã}{{\~a}}1 {õ}{{\~o}}1
}
\lstset{style=shell}

\setlength{\parindent}{1.25cm}

% =============================================================================
% Documento
% =============================================================================
\begin{document}

% --------------------------------------------------------------------------
% Capa
% --------------------------------------------------------------------------
\begin{titlepage}
    \centering
    \vspace*{1cm}

    {\large NOME DA INSTITUIÇÃO DE ENSINO}\\[0.3cm]
    {\large CURSO DE [NOME DO CURSO]}\\[0.3cm]
    {\large DISCIPLINA DE SISTEMAS OPERACIONAIS}

    \vspace{3cm}

    {\Huge \bfseries Linux Mint em Ambiente Docker}\\[0.5cm]
    {\Large Estrutura do Sistema Operacional, Kernel e Automação via Shell Script}

    \vspace{3cm}

    {\large Avaliação AV2}

    \vspace{2cm}

    {\large \textbf{Equipe:}}\\[0.3cm]
    \begin{tabular}{c}
        Nome do Integrante 1 \\
        Nome do Integrante 2 \\
        Nome do Integrante 3 \\
        Nome do Integrante 4 \\
        Nome do Integrante 5 \\
    \end{tabular}

    \vfill

    {\large Maceió/AL}\\
    {\large 2026}

\end{titlepage}

% --------------------------------------------------------------------------
% Sumário
% --------------------------------------------------------------------------
\tableofcontents
\newpage

% --------------------------------------------------------------------------
% 1. Introdução
% --------------------------------------------------------------------------
\section{Introdução}

Este relatório foi desenvolvido como parte da Avaliação AV2 da disciplina de
Sistemas Operacionais, cujo objetivo é o estudo prático e teórico de um
sistema operacional executado em ambiente virtualizado por containers. O
sistema operacional escolhido pela equipe foi o \textbf{Linux Mint}, uma
distribuição Linux baseada no Ubuntu/Debian, amplamente utilizada em desktops
por sua estabilidade, leveza e familiaridade de uso.

O trabalho contempla os seguintes pontos, conforme especificado nas regras da
atividade:

\begin{itemize}
    \item Instalação do sistema operacional escolhido em um ambiente
          containerizado (Docker);
    \item Apresentação da estrutura geral do sistema operacional;
    \item Explicação sobre o funcionamento do Kernel;
    \item Demonstração prática utilizando uma linguagem de script
          (Shell Script);
    \item Versionamento de todo o material produzido em um repositório
          do GitHub;
    \item Elaboração deste relatório em \LaTeX.
\end{itemize}

Nas seções a seguir, será apresentada a justificativa para o uso do Docker,
uma visão geral do Linux Mint, a estrutura de diretórios do sistema
operacional (FHS), o funcionamento do Kernel Linux, os scripts desenvolvidos
em Shell Script e, por fim, a estratégia de versionamento utilizada.


% --------------------------------------------------------------------------
% 2. Ambiente do Projeto: Docker
% --------------------------------------------------------------------------
\section{Ambiente do Projeto: Docker}

\subsection{O que é o Docker}

O Docker é uma plataforma de virtualização baseada em \textbf{containers}.
Diferentemente de uma máquina virtual tradicional, que emula um hardware
completo e exige um sistema operacional convidado completo (com seu próprio
kernel), um container compartilha o kernel do sistema hospedeiro e isola
apenas o espaço de usuário (bibliotecas, binários e arquivos de
configuração). Isso resulta em imagens mais leves, criação quase instantânea
e menor consumo de recursos quando comparado a uma máquina virtual completa.

A equipe optou pelo uso do Docker em vez de uma máquina virtual completa pelos
seguintes motivos:

\begin{itemize}
    \item \textbf{Portabilidade}: o ambiente pode ser reproduzido de forma
          idêntica em qualquer computador que tenha o Docker instalado,
          independentemente do sistema operacional do hospedeiro;
    \item \textbf{Reprodutibilidade}: o \texttt{Dockerfile} documenta de
          forma explícita e versionável todos os passos necessários para
          montar o ambiente;
    \item \textbf{Leveza}: a imagem ocupa menos espaço e inicializa em
          poucos segundos, facilitando testes repetidos durante o
          desenvolvimento dos scripts.
\end{itemize}

\subsection{Configuração do Linux Mint no Docker}

Para executar o Linux Mint em container, foi utilizada uma imagem baseada no
Linux Mint 21.3 ("Virginia"), que por sua vez é construída sobre o Ubuntu 22.04
LTS. A partir dessa imagem base, foi criado um \texttt{Dockerfile}
personalizado que:

\begin{enumerate}
    \item Instala ferramentas auxiliares utilizadas pelos scripts da equipe
          (\texttt{htop}, \texttt{tree}, \texttt{neofetch},
          \texttt{net-tools}, entre outras);
    \item Copia os scripts Shell desenvolvidos pela equipe para dentro do
          container, no diretório \texttt{/projeto/scripts};
    \item Define como comando padrão (\texttt{CMD}) a execução do menu
          interativo \texttt{00\_menu.sh}, que dá acesso a todas as
          demonstrações realizadas.
\end{enumerate}

O trecho a seguir apresenta o \texttt{Dockerfile} utilizado:

\begin{lstlisting}[caption={Dockerfile do ambiente Linux Mint}]
FROM linuxmintd/mint21.3-amd64:latest

ENV DEBIAN_FRONTEND=noninteractive

RUN apt-get update && apt-get install -y \
        htop tree neofetch net-tools procps \
        nano less man-db kmod \
    && apt-get clean && rm -rf /var/lib/apt/lists/*

WORKDIR /projeto

COPY scripts/ /projeto/scripts/

RUN chmod +x /projeto/scripts/*.sh

CMD ["/bin/bash", "/projeto/scripts/00_menu.sh"]
\end{lstlisting}

Para construir e executar o container, são utilizados os comandos:

\begin{lstlisting}[language=bash, caption={Construção e execução do container}]
# Construir a imagem
docker build -t linuxmint-av2 -f docker/Dockerfile .

# Executar o container interativamente
docker run -it --rm linuxmint-av2
\end{lstlisting}

Alternativamente, a equipe disponibilizou um arquivo
\texttt{docker-compose.yml}, permitindo a execução com um único comando:

\begin{lstlisting}[language=bash, caption={Execução via Docker Compose}]
docker compose -f docker/docker-compose.yml run --rm linuxmint-av2
\end{lstlisting}


% --------------------------------------------------------------------------
% 3. O Sistema Operacional Linux Mint
% --------------------------------------------------------------------------
\section{O Sistema Operacional Linux Mint}

\subsection{Visão geral}

O \textbf{Linux Mint} é uma distribuição Linux de código aberto, lançada em
2006, derivada do Ubuntu (e, por consequência, do Debian). Seu principal
objetivo é oferecer um sistema operacional elegante, moderno, confortável e
de fácil utilização, sendo uma das distribuições mais populares para
computadores desktop. Entre suas principais características, destacam-se:

\begin{itemize}
    \item Utiliza o ambiente de área de trabalho \textbf{Cinnamon} (em sua
          edição principal), além de oferecer variações com os ambientes
          MATE e Xfce;
    \item Mantém compatibilidade total com o ecossistema de pacotes
          \texttt{.deb} e com o gerenciador de pacotes \texttt{APT},
          herdado do Ubuntu/Debian;
    \item Inclui, por padrão, codecs multimídia e drivers proprietários,
          facilitando o uso imediato após a instalação;
    \item É amplamente utilizado como alternativa ao Windows, por usuários
          que buscam uma interface familiar com menor custo de
          aprendizado.
\end{itemize}

\subsection{Estrutura do Sistema Operacional}

De forma geral, qualquer sistema operacional do tipo Linux pode ser
representado em camadas, conforme ilustrado na Figura \ref{fig:camadas}:

\begin{enumerate}
    \item \textbf{Hardware}: processador, memória, discos, placas de rede
          e demais dispositivos físicos;
    \item \textbf{Kernel}: núcleo do sistema operacional, responsável por
          gerenciar diretamente o hardware e oferecer uma interface
          padronizada (chamadas de sistema) para os programas;
    \item \textbf{Shell e bibliotecas do sistema}: interpretadores de
          comandos (como o Bash) e bibliotecas (como a \texttt{glibc}) que
          fazem a ponte entre as aplicações e o kernel;
    \item \textbf{Aplicações}: programas utilizados pelo usuário final
          (navegadores, editores de texto, ferramentas administrativas
          etc.).
\end{enumerate}

\begin{figure}[h]
    \centering
    \fbox{\begin{minipage}{0.7\textwidth}
        \centering
        \vspace{0.3cm}
        \textbf{Aplicações} (LibreOffice, Firefox, Terminal, ...) \\[0.4cm]
        $\big\downarrow$ chamadas de biblioteca $\big\uparrow$ \\[0.4cm]
        \textbf{Shell / Bibliotecas} (Bash, glibc) \\[0.4cm]
        $\big\downarrow$ chamadas de sistema (syscalls) $\big\uparrow$ \\[0.4cm]
        \textbf{Kernel Linux} (processos, memória, arquivos, rede,
        drivers) \\[0.4cm]
        $\big\downarrow$ instruções de hardware $\big\uparrow$ \\[0.4cm]
        \textbf{Hardware} (CPU, RAM, Disco, Rede)
        \vspace{0.3cm}
    \end{minipage}}
    \caption{Camadas de um sistema operacional Linux}
    \label{fig:camadas}
\end{figure}

\subsubsection{Hierarquia de diretórios (FHS)}

A organização do sistema de arquivos do Linux Mint segue o padrão
\textbf{FHS} (\textit{Filesystem Hierarchy Standard}), comum à maioria das
distribuições Linux. A Tabela \ref{tab:fhs} resume os diretórios mais
importantes encontrados na raiz (\texttt{/}) do sistema.

\begin{table}[h]
    \centering
    \begin{tabular}{|l|p{10cm}|}
        \hline
        \textbf{Diretório} & \textbf{Descrição} \\ \hline
        \texttt{/bin}  & Binários essenciais (comandos básicos como
                          \texttt{ls}, \texttt{cp}, \texttt{cat}). \\ \hline
        \texttt{/boot} & Arquivos do kernel (\texttt{vmlinuz}), initramfs
                          e do gerenciador de boot (GRUB). \\ \hline
        \texttt{/dev}  & Arquivos especiais que representam dispositivos
                          de hardware. \\ \hline
        \texttt{/etc}  & Arquivos de configuração do sistema e dos
                          aplicativos. \\ \hline
        \texttt{/home} & Diretórios pessoais dos usuários. \\ \hline
        \texttt{/lib}  & Bibliotecas compartilhadas usadas pelos
                          binários do sistema. \\ \hline
        \texttt{/proc} & Sistema de arquivos virtual com informações do
                          kernel e processos em execução. \\ \hline
        \texttt{/sys}  & Sistema de arquivos virtual com informações
                          sobre dispositivos e drivers. \\ \hline
        \texttt{/tmp}  & Arquivos temporários. \\ \hline
        \texttt{/usr}  & Programas, bibliotecas e documentação para os
                          usuários (maior parte do sistema). \\ \hline
        \texttt{/var}  & Dados variáveis: logs, caches, bancos de dados
                          de pacotes. \\ \hline
    \end{tabular}
    \caption{Principais diretórios do FHS no Linux Mint}
    \label{tab:fhs}
\end{table}

Essa estrutura pode ser explorada de forma prática por meio do script
\texttt{02\-\_es\-tru\-tu\-ra\-\_ar\-qui\-vos.sh}, descrito na Seção
\ref{sec:scripts}.


% --------------------------------------------------------------------------
% 4. O Kernel Linux
% --------------------------------------------------------------------------
\section{O Kernel Linux}

\subsection{Definição e tipo de kernel}

O \textbf{kernel} é o componente central de um sistema operacional. Ele atua
como intermediário entre o hardware e os programas em execução, sendo
responsável por controlar o acesso aos recursos físicos da máquina (CPU,
memória, dispositivos de entrada e saída) de forma segura e organizada.

Existem diferentes arquiteturas de kernel, sendo as principais:

\begin{itemize}
    \item \textbf{Microkernel}: apenas funcionalidades mínimas (comunicação
          entre processos, gerenciamento básico de memória e
          escalonamento) ficam no espaço do kernel; drivers e serviços
          rodam em espaço de usuário, separados em processos isolados;
    \item \textbf{Kernel monolítico}: todos os serviços principais do
          sistema operacional (gerenciamento de processos, memória,
          sistema de arquivos, drivers, rede) executam no mesmo espaço de
          endereçamento privilegiado (kernel space).
\end{itemize}

O \textbf{Linux} adota uma abordagem \textbf{monolítica modular}: o núcleo é
monolítico, mas grande parte de seus componentes (especialmente drivers de
dispositivos) pode ser compilada como \textbf{módulos} (arquivos
\texttt{.ko}), que são carregados e descarregados dinamicamente em tempo de
execução, sem necessidade de reiniciar o sistema ou recompilar o kernel.

\subsection{Principais funções do Kernel}

\begin{enumerate}
    \item \textbf{Gerenciamento de processos}: o kernel cria, escalona,
          suspende e finaliza processos. O componente responsável por
          decidir qual processo utilizará a CPU em cada instante é o
          \textit{scheduler} (escalonador);
    \item \textbf{Gerenciamento de memória}: controla a alocação de memória
          física e virtual para cada processo, implementa a memória
          \textit{swap} e garante o isolamento entre os espaços de memória
          dos processos;
    \item \textbf{Gerenciamento de dispositivos (drivers)}: comunica-se com
          o hardware por meio de drivers, que no Linux frequentemente são
          implementados como módulos do kernel;
    \item \textbf{Sistema de arquivos}: implementa o suporte a diferentes
          formatos de sistemas de arquivos (ext4, btrfs, NTFS, etc.) e
          fornece chamadas de sistema (\textit{syscalls}) como
          \texttt{open}, \texttt{read}, \texttt{write} e \texttt{close};
    \item \textbf{Comunicação entre processos e rede}: implementa
          mecanismos de IPC (\textit{pipes}, sinais, memória compartilhada)
          e toda a pilha de protocolos de rede (TCP/IP).
\end{enumerate}

\subsection{Observando o Kernel na prática}

Dentro do container Linux Mint, é possível observar informações fornecidas
pelo kernel por meio de comandos simples e dos sistemas de arquivos virtuais
\texttt{/proc} e \texttt{/sys}. Alguns exemplos utilizados no script
\texttt{03\_kernel\_info.sh}:

\begin{lstlisting}[language=bash, caption={Comandos para inspecionar o Kernel}]
# Versão do kernel em execução
uname -a

# Detalhes da versão do kernel (compilador, data)
cat /proc/version

# Módulos do kernel carregados atualmente
lsmod

# Parâmetros do kernel em tempo de execução
sysctl kernel.ostype kernel.osrelease kernel.version

# Dispositivos gerenciados pelo kernel
ls /sys/class
\end{lstlisting}

Esses comandos demonstram, na prática, que o kernel está constantemente
expondo informações sobre o estado do hardware e dos processos através de
arquivos virtuais, permitindo que ferramentas de espaço de usuário (como
\texttt{ps}, \texttt{top} e \texttt{free}) leiam e apresentem esses dados de
forma amigável.


% --------------------------------------------------------------------------
% 5. Linguagem utilizada: Shell Script
% --------------------------------------------------------------------------
\section{Linguagem de Programação Utilizada: Shell Script}
\label{sec:scripts}

\subsection{Justificativa}

A linguagem escolhida pela equipe para a automação das demonstrações foi o
\textbf{Shell Script}, na variante \textbf{Bash} (\textit{Bourne Again
Shell}), o interpretador de comandos padrão do Linux Mint. A escolha se
justifica por:

\begin{itemize}
    \item Ser a linguagem nativa de interação com o sistema operacional,
          permitindo acesso direto a comandos, arquivos de configuração e
          aos sistemas de arquivos virtuais \texttt{/proc} e \texttt{/sys};
    \item Não exigir instalação de ferramentas adicionais, já que o Bash
          está presente em qualquer instalação padrão do Linux Mint
          (e, portanto, na imagem Docker utilizada);
    \item Ser amplamente utilizada para automação de tarefas
          administrativas, sendo um conteúdo diretamente relacionado à
          disciplina de Sistemas Operacionais.
\end{itemize}

\subsection{Scripts desenvolvidos}

Foram desenvolvidos seis scripts, organizados no diretório
\texttt{scripts/} do repositório, descritos na Tabela \ref{tab:scripts}.

\begin{table}[h]
    \centering
    \begin{tabular}{|l|p{9.5cm}|}
        \hline
        \textbf{Script} & \textbf{Descrição} \\ \hline
        \texttt{00\_menu.sh} &
            Menu interativo, ponto de entrada do container, que permite
            escolher e executar as demais demonstrações. \\ \hline
        \texttt{01\_info\_sistema.sh} &
            Exibe informações gerais do sistema: distribuição,
            \textit{hostname}, tempo de atividade e arquitetura. \\ \hline
        \texttt{02\_estrutura\_arquivos.sh} &
            Lista e descreve a estrutura de diretórios do sistema
            operacional (FHS). \\ \hline
        \texttt{03\_kernel\_info.sh} &
            Exibe informações sobre o kernel: versão, módulos carregados,
            parâmetros de execução e explica suas funções. \\ \hline
        \texttt{04\_processos\_memoria.sh} &
            Demonstra o gerenciamento de processos e memória realizado
            pelo kernel, utilizando \texttt{ps}, \texttt{free} e
            \texttt{pstree}. \\ \hline
        \texttt{05\_pacotes.sh} &
            Demonstra o funcionamento do gerenciador de pacotes
            APT/dpkg, utilizado pelo Linux Mint. \\ \hline
    \end{tabular}
    \caption{Scripts desenvolvidos pela equipe}
    \label{tab:scripts}
\end{table}

A título de exemplo, o trecho a seguir apresenta o script
\texttt{02\-\_es\-tru\-tu\-ra\-\_ar\-qui\-vos.sh}, responsável por listar e descrever os
diretórios da raiz do sistema:

\begin{lstlisting}[language=bash, caption={Trecho do script 02\_estrutura\_arquivos.sh}]
#!/bin/bash
echo "Conteúdo da raiz '/':"
ls -lh /

echo "Descrição de cada diretório principal:"
cat <<'TEXTO'
/bin    -> Binários essenciais usados por todos os usuários
/boot   -> Arquivos necessários para o boot (kernel, GRUB)
/etc    -> Arquivos de configuração do sistema
/home   -> Diretórios pessoais dos usuários
/proc   -> Informações do kernel e processos (sistema virtual)
/usr    -> Programas e bibliotecas para os usuários
/var    -> Dados variáveis: logs, caches, etc.
TEXTO
\end{lstlisting}

Todos os scripts foram testados dentro do container Linux Mint e podem ser
executados individualmente ou através do menu principal
(\texttt{00\_menu.sh}), conforme demonstrado durante a apresentação.


% --------------------------------------------------------------------------
% 6. Versionamento com Git e GitHub
% --------------------------------------------------------------------------
\section{Versionamento com Git e GitHub}

Todo o material produzido pela equipe (Dockerfile, scripts, relatório em
\LaTeX{} e demais arquivos de apoio) foi versionado utilizando o sistema de
controle de versão \textbf{Git} e hospedado em um repositório do
\textbf{GitHub}.

A organização do repositório segue a estrutura apresentada a seguir:

\begin{lstlisting}[caption={Estrutura do repositório no GitHub}]
projeto-av2-linux-mint/
|-- README.md
|-- .gitignore
|-- docker/
|   |-- Dockerfile
|   |-- docker-compose.yml
|   `-- .dockerignore
|-- scripts/
|   |-- 00_menu.sh
|   |-- 01_info_sistema.sh
|   |-- 02_estrutura_arquivos.sh
|   |-- 03_kernel_info.sh
|   |-- 04_processos_memoria.sh
|   `-- 05_pacotes.sh
`-- relatorio/
    |-- relatorio.tex
    `-- relatorio.pdf
\end{lstlisting}

O fluxo de trabalho adotado pela equipe seguiu, de forma simplificada, os
seguintes passos:

\begin{enumerate}
    \item Criação do repositório no GitHub e clonagem local
          (\texttt{git clone});
    \item Desenvolvimento incremental do \texttt{Dockerfile} e dos
          scripts, com commits descritivos para cada funcionalidade
          adicionada;
    \item Testes do ambiente Docker a cada alteração relevante;
    \item Escrita e atualização do relatório em \LaTeX{} em paralelo ao
          desenvolvimento técnico;
    \item Envio (\textit{push}) das alterações para o repositório remoto,
          mantendo o histórico de versões disponível para consulta.
\end{enumerate}

O link do repositório utilizado pela equipe é:
\url{https://github.com/USUARIO/projeto-av2-linux-mint}
\textit{(substituir pelo endereço real do repositório da equipe)}.


% --------------------------------------------------------------------------
% 7. Conclusão
% --------------------------------------------------------------------------
\section{Conclusão}

O desenvolvimento deste projeto permitiu à equipe consolidar, de forma
prática, diversos conceitos fundamentais da disciplina de Sistemas
Operacionais. A utilização do Docker possibilitou a criação de um ambiente
isolado, portátil e facilmente reproduzível, contendo o Linux Mint, sem a
necessidade de uma instalação tradicional em máquina física ou virtual
completa.

Através da exploração da estrutura de diretórios (FHS) e dos mecanismos
oferecidos pelo Kernel Linux (gerenciamento de processos, memória,
dispositivos e arquivos), foi possível compreender, na prática, como as
camadas de hardware, kernel, shell e aplicações se relacionam para formar um
sistema operacional funcional.

Por fim, o uso de Shell Script para automatizar as demonstrações reforçou a
importância dessa linguagem como ferramenta de administração de sistemas
Linux, enquanto o versionamento com Git e GitHub garantiu organização,
rastreabilidade e colaboração eficiente entre os membros da equipe durante
todo o desenvolvimento do projeto.


% --------------------------------------------------------------------------
% 8. Referências
% --------------------------------------------------------------------------
\section{Referências}

\sloppy
\begin{enumerate}
    \item LINUX MINT TEAM. \textit{Linux Mint Documentation}. Disponível em:
          \url{https://linuxmint.com}.
    \item THE LINUX KERNEL ORGANIZATION. \textit{The Linux Kernel
          documentation}. Disponível em: \url{https://kernel.org}.
    \item DOCKER INC. \textit{Docker Documentation}. Disponível em:
          \url{https://docs.docker.com}.
    \item FILESYSTEM HIERARCHY STANDARD GROUP. \textit{Filesystem
          Hierarchy Standard (FHS)}. Disponível em:
          \url{https://refspecs.linuxfoundation.org/fhs.shtml}.
    \item TANENBAUM, A. S.; BOS, H. \textit{Modern Operating Systems}.
          4. ed. Pearson, 2014.
\end{enumerate}

\end{document}
